\section{Bilinear pairings on elliptic curves}\label{pair}

Before we are ready for KZG, there is one more piece of elliptic curve
math that we need.

Recall that the map $[ \bullet ]:\FF_{q} \rightarrow E$ is linear,
meaning that
$[a + b] = [a] + [b]$, and $[na] = n [a]$.
But as written we cannot do ``armored multiplication'':

\begin{assumption}{}{}
As far as we know, given $[a]$ and $[b]$, one cannot easily compute $[a b]$.
\end{assumption}

On the other hand, it \textbf{does} turn out that we know a way to
\textbf{verify} a claimed answer on certain curves. That is:

\begin{proposition}{BN254 is pairing-friendly}{}
  On the curve BN254, given three points $[a]$,
  $[b]$, and $[c]$ on the curve, one can verify whether $ab = c$.
\end{proposition}

The technique needed is that one wants to construct a nondegenerate
bilinear function
\[ \pair \colon E \times E \to \ZZ/N\ZZ \]
for some large integer $N$.
We think this should be called a \emph{bilinear pairing},
but for some reason everyone just says \alert{pairing} instead.
A curve is called \alert{pairing-friendly} if
this pairing can be computed reasonably quickly
(e.g., BN254 is pairing-friendly, but Curve25519 is not).

This construction actually uses some really deep number theory
(heavier than all the math in Section \ref{ec})
that is well beyond the scope of this modest book.
Fortunately, we will not need the details of how it works;
but we will comment briefly in Section \ref{pairing-friendly} on what curves it
can be done on.
And this pairing algorithm needs to be worked out just
once for the curve $E$;
and then anyone in the world can use the published curve for their protocol.

Going a little more generally, the four-number equation
\[ \pair([m],[n]) = \pair([m'],[n']) \]
will be true whenever $mn = m'n'$, because both sides will equal
$mn \pair([1], [1])$.
So this gives us a way to \textbf{verify} two-by-two multiplication.

\begin{remark}{}{}
The last sentence is worth bearing in mind: In all the protocols we will see,
the pairing is only used by the verifier Victor, never by the prover Peggy.
\end{remark}

\begin{remark}{}{}
On the other hand, we currently do not seem to know a good
way to do \emph{multilinear} pairings.
For example, we do not know a good trilinear map
$E \times E \times E \rightarrow \ZZ / N \ZZ$
that would allow us to compare $[a b c]$, $[a]$, $[b]$, $[c]$
(without knowing one of $[a b]$, $[b c]$, $[c a]$).
\end{remark}

\subsection{Verifying more complicated claims}\label{pair-verify}

\begin{example}{}{}
Suppose Peggy wants to convince Victor that $y = x^{3} + 2$, where
Peggy has sent Victor elliptic curve points $[x]$ and
$[y]$. To do this, Peggy additionally sends to Victor
$[x^2]$ and
$[x^3]$.

Given $[x]$, $[x^2]$,
$[x^3]$, and $[y]$, Victor
verifies that:

\begin{itemize}
\item
  $\pair([x^{2}],[1]) = \pair([x],[x])$
\item
  $\pair([x^{3}],[1]) = \pair([x^2],[x])$
\item
  $[y] = [x^3] + 2[1]$.
\end{itemize}

The process of verifying this sort of identity is quite general: The
prover sends intermediate values as needed so that the verifier can
verify the claim using only pairings and linearity.
\end{example}

\subsection{So which curves are pairing-friendly?}\label{pairing-friendly}

If we chose $E$ to be BN254, the following property holds:

\begin{proposition}{}{}
For $(p,q)$ as in BN254, the smallest integer $k$ such that $q$
divides $p^{k} - 1$ is $k = 12$.
\end{proposition}

This integer $k$ is called the \alert{embedding degree}. This section
is an aside explaining how the embedding degree affects pairing.

The pairing function $\pair(a,b)$ takes as input two
points $a,b \in E$ on the elliptic curve, and spits out a value
$\pair(a,b) \in \FF_{p^{k}}^{\times}$ -- in
other words, a nonzero element of the finite field of order $p^{k}$
(where $k$ is the embedding degree we just defined). In fact, this
element will always be a $q$-th root of unity in
$\FF_{p^{k}}$, and it will satisfy
$\pair([m],[n]) = \zeta^{mn}$,
where $\zeta$ is some fixed $q$-th root of unity. The construction
of the pairing is based on the Weil pairing\footurl{https://en.wikipedia.org/wiki/Weil\_pairing}
in algebraic geometry. How to compute these pairings is well beyond the
scope of these notes; the raw definition is quite abstract, and a lot of
work has gone into computing the pairings efficiently. (For more
details, see these notes.\footurl{https://crypto.stanford.edu/pbc/notes/ep/pairing.html})

The difficulty of computing these pairings is determined by the size of
$k$: The values $\pair(a,b)$ will be elements of a
field of size $p^{k}$, so they will require $256k$ bits even to
store. For a curve to be ``pairing-friendly'' -- in order to be able to
do pairing-based cryptography on it -- we need the value of $k$ to be
pretty small.
